\begin{circuitikz}[scale=0.9, every node/.style={transform shape}, font=\sffamily\small]
% ---------------- the switch in the middle
\draw[rounded corners=2pt, fill=extcol] (5.0,2.2) rectangle (9.0,3.4);
\node[font=\sffamily\bfseries\small] at (7.0,3.05) {LAN switch};
\node[font=\sffamily\scriptsize] at (7.0,2.6) {or any access point on the bench};
% ---------------- the operator
\draw[rounded corners=2pt, fill=usrcol] (0.2,2.2) rectangle (3.6,3.4);
\node[font=\sffamily\bfseries\small] at (1.9,3.05) {operator Pi};
\node[font=\sffamily\scriptsize] at (1.9,2.6) {holds the private key};
\draw[thick] (3.6,2.8) -- (5.0,2.8);
% ---------------- the four hosts, hung below the switch
\foreach \x/\name/\sub/\port in {0.7/pi4/{Raspberry Pi 4}/5.6, 4.1/pi3bp/{Raspberry Pi 3B+}/6.4,
                                 7.5/pi3/{Raspberry Pi 3}/7.6, 10.9/air/{NanoPi NEO Air}/8.4}{
  \draw[rounded corners=2pt, fill=hwcol] (\x,-0.4) rectangle (\x+3.0,0.9);
  \node[font=\sffamily\bfseries\small] at (\x+1.5,0.55) {\name};
  \node[font=\sffamily\scriptsize] at (\x+1.5,0.1) {\sub};
  \draw (\x+1.5,0.9) -- (\x+1.5,1.6) -- (\port,1.6) -- (\port,2.2);
}
\node[anchor=west, font=\sffamily\scriptsize] at (0.2,1.85) {Ethernet, or Wi-Fi for the NanoPi};
% ---------------- what each link carries
\node[anchor=west, font=\sffamily\scriptsize, align=left] at (0.2,-1.3)
  {Every link carries one thing: an SSH session that runs read-only commands and returns text.};
\node[anchor=west, font=\sffamily\scriptsize, text=red!70!black, align=left] at (0.2,-1.9)
  {The 40-pin headers are read in this lab and never changed. No modem is reconfigured.};
\node[anchor=west, font=\sffamily\scriptsize, text=black!60, align=left] at (0.2,-2.5)
  {The NanoPi joined its access point in P09. Its header bus is 0, not 1.};
\end{circuitikz}
